\documentclass[]{usiinfbachelorproject}

\captionsetup{labelfont={bf}}

\author{Andrea Vicari}

\title{Smart--IVC}
\subtitle{Cities Becoming Alive}
\versiondate{\today}
\begin{committee}
%With more than 1 advisor an error is raised...: only 1 advisor is allowed!
\advisor[Universit\`a della Svizzera Italiana, Switzerland]{Prof.}{Michele}{Lanza}
%You can comment out  these lines if you don't have any assistant
\assistant[Universit\`a della Svizzera Italiana, Switzerland]{Prof. Dr.}{Andrea}{Mocci}
\end{committee}


\abstract {
%The purpose of this Bachelor Project is to create a web application that provides a 3D-visualization a city and let the user interact with the entities inside it.\\
Web users have to aggregate different data from various sources whenever they are looking for some information on the Internet. Think about the student who is looking for a rented room near her university: she firstly uses a specific website to find the advertisement of a room for rent, she then looks for the address on another website that provides a map service to see if the house is located where she desires.\\

No such service exists that provides a unique environment in which the user can both visualize a city and interact with its elements. The only technologies available are either not exhaustive or too complex to use.\\

This thesis introduces Smart--IVC a web application that provides an intuitive interface and prevents the user from jumping from website to website. Through the form of a 3D-environment, this application provides an interactive visualisation of cities in which the user can directly “communicate” with the elements, executing queries on them. After having clicked on a building in the map, the user is able to get information (coordinates, address, floors etc.) about that construction and also find out what are the various relations between that specific building and the other entities in the city.\\

Smart--IVC is an application accessible by everyone and that aims to enhance the city visualisation, getting closer to the user needs.
}


\begin{document}
\maketitle

%%%%%%%%%%%%%%%%%%%%%%%%%
\section{Introduction} \label{introduction}
%%%%%%%%%%%%%%%%%%%%%%%%%

%%%%%%%%%%%%%%%%%%%%%%%%%
\section{State of the Art} \label{stateOfTheArt}
%%%%%%%%%%%%%%%%%%%%%%%%%

%%%%%%%%%%%%%%%%%%%%%%%%%
\section{Project requirements and analysis} \label{projectRequirementsAndAnalysis}
%%%%%%%%%%%%%%%%%%%%%%%%%

%%%%%%%%%%%%%%%%%%%%%%%%%
\section{Project design} \label{projectDesign}
%%%%%%%%%%%%%%%%%%%%%%%%%

%%%%%%%%%%%%%%%%%%%%%%%%%
\section{Implementation issues} \label{implementationIssues}
%%%%%%%%%%%%%%%%%%%%%%%%%

%%%%%%%%%%%%%%%%%%%%%%%%%
\section{Tests (methodology, results, comments)} \label{tests}
%%%%%%%%%%%%%%%%%%%%%%%%%

%%%%%%%%%%%%%%%%%%%%%%%%%
\section{Conclusions and future work or possible developments} \label{conclusions}
%%%%%%%%%%%%%%%%%%%%%%%%%


%\begin{figure} [h]
%\centering
%\includegraphics[width=0.5\textwidth]{logo-info.pdf}
%\caption{Caption of the figure}
%\label{fig:USILogo}
%\end{figure}

%\begin{table}[h]
%\centering
%\scalebox {0.8} {
%\begin{normalsize}\begin{tabular}{l|lll}
%\textbf{Col 1} & \textbf{Col 2} & \textbf{Col 3} & \textbf{Col 4}\\
%\hline
%1 & 2 & 3 & Goofy\\
%4 & 5 & 6 & Mickey
%\end{tabular}
%\end{normalsize}
%}
%\caption{Caption of the table}
%\label{tab:numbers}
%\end{table}

%\cite{Stru1899a}
%%%%%
\bibliographystyle{abbrv}
\bibliography{references}

\end{document}